\documentclass[12pt]{article}
\usepackage[T1]{fontenc}
\usepackage[utf8]{inputenc}
\usepackage{lmodern}
\usepackage{textcomp}
\usepackage{enumitem}
\usepackage{geometry}
\geometry{margin=1in}
\begin{document}
\begin{center}
Answer Key - Biology - Undergraduate Level Assessment 13 \
\small (Answers are ordered exactly as in the question paper.)
\end{center}

\section*{Multiple Choice}
\begin{enumerate}[label=\arabic*.]
\item \textbf{Answer:} A
\\ \textit{Options:} A) Human sperm cannot develop at the high temperature found within the body cavity.; B) The testes require the lower pressure environment of the scrotum for sperm maturation; C) Testes must be in the scrotum to properly attach to the vas deferens; D) The scrotum provides the necessary nutrients for sperm development; E) Gravity is required for the testes to function properly
\\ \textit{Note:} The correct answer is based on the fact that human sperm require a cooler environment for proper development, which is achieved in the scrotum, as the higher temperatures in the body cavity can adversely affect sperm production and viability.
\item \textbf{Answer:} A
\\ \textit{Options:} A) Carbon dioxide bubbles forming in the blood causing acidosis; B) Rapid increase in barometric pressure causing the bends; C) Formation of helium bubbles in the bloodstream due to rapid ascent; D) Increased salinity in the water leading to dehydration and joint pain; E) Oxygen bubbles forming in the blood leading to hypoxia
\\ \textit{Note:} The bends, or decompression sickness, occurs when dissolved gases, primarily nitrogen, come out of solution to form bubbles in the bloodstream during rapid ascent, leading to various complications including acidosis due to carbon dioxide accumulation from tissue hypoxia and bubble-related damage.
\item \textbf{Answer:} A
\\ \textit{Options:} A) Protects the cell from foreign invaders; B) Acts as a storage for genetic material; C) Regulates the pH level within the cell; D) Transports oxygen throughout the cell; E) Removes nitrogenous wastes
\\ \textit{Note:} The contractile vacuole plays a crucial role in maintaining osmotic balance by expelling excess water and waste, which helps protect the cell from potential damage and foreign invaders.
\item \textbf{Answer:} A
\\ \textit{Options:} A) The human genome has fully adapted to current environmental pressures; B) Modern medicine has eliminated the effects of genetic drift; C) The rate of genetic drift is directly proportional to technological advancements; D) Humans have reached a genetic equilibrium where drift is no longer possible; E) Humans have developed immunity to genetic drift
\\ \textit{Note:} The correct answer is based on the understanding that as the human genome has become more stable and adapted to current environmental pressures, the influence of random changes in allele frequencies due to genetic drift has diminished compared to earlier periods in human history when populations were more isolated and subject to varying selective pressures.
\item \textbf{Answer:} B
\\ \textit{Options:} A) The absence of a Y chromosome determines maleness in mammals; B) The presence of a Y chromosome usually determines maleness in mammals; C) The ratio of two X chromosomes to one Y chromosome determines maleness; D) The presence of an X chromosome determines maleness in mammals; E) The presence of an X chromosome in the absence of a Y chromosome determines maleness
\\ \textit{Note:} The presence of a Y chromosome typically leads to the development of male characteristics in mammals due to the SRY gene it carries, which initiates the pathway for male sex determination.
\end{enumerate}

\section*{True / False}
\begin{enumerate}[label=\arabic*.]
\setcounter{enumi}{5}
\item \textbf{Answer:} False
\\ \textit{Options:} A) True; B) False
\\ \textit{Note:} Oxidative phosphorylation primarily produces ATP through the electron transport chain and chemiosmosis, while NADH and FADH 2 are produced during earlier stages of cellular respiration, such as glycolysis and the citric acid cycle, rather than during oxidative phosphorylation itself.
\item \textbf{Answer:} True
\\ \textit{Options:} A) True; B) False
\\ \textit{Note:} The statement is true because the gal^+ phenotype in a gal^+/gal^- cell produced by abortive transduction is defective in galactose metabolism, leading to slower growth and smaller colonies when grown on galactose as the only carbon source.
\item \textbf{Answer:} False
\\ \textit{Options:} A) True; B) False
\\ \textit{Note:} Proteins destined for the lysosome are typically targeted by a single signal, known as the mannose-6-phosphate (M$_{6}$P) marker, which is sufficient for their recognition and transport to the lysosome.
\item \textbf{Answer:} False
\\ \textit{Options:} A) True; B) False
\\ \textit{Note:} The stages of a behavioral act typically include cue, motivation, and action, rather than desire and fulfillment, making the statement inaccurate.
\item \textbf{Answer:} True
\\ \textit{Options:} A) True; B) False
\\ \textit{Note:} The acrosome reaction allows a single sperm to penetrate the egg's protective layers, while the cortical reaction modifies the egg's membrane to prevent additional sperm entry, thereby ensuring that only one sperm fertilizes the egg and preventing the formation of twin zygotes.
\end{enumerate}

\section*{Long Form Response}
\begin{enumerate}[label=\arabic*.]
\setcounter{enumi}{10}
\item \textbf{Answer:} The gene frequency of 0.07 for widow's peak, a dominant trait, indicates that the recessive allele frequency is 0.93 (1 - 0.07). Using the Hardy-Weinberg principle, the expected frequencies of the genotypes can be calculated as follows: homozygous dominant ($p^2$) is 0.$07^2$ = 0.0049, heterozygous (2 pq) is 2(0.07)(0.93) = 0.1302, and homozygous recessive ($q^2$) is 0.$93^2$ = 0.8649. This suggests that in a large population, approximately 0.49\% would be homozygous dominant, 13.02\% would be heterozygous, and 86.49\% would be homozygous recessive, indicating that the widow's peak trait is likely to be present in a small fraction of the population, with the majority being non-expressive of the trait. This distribution has implications for genetic diversity and the potential for trait expression within the population.
\\ \textit{Note:} correct because it accurately applies the Hardy-Weinberg principle to calculate the expected frequencies of homozygous dominant, heterozygous, and homozygous recessive individuals based on the gene frequency of the dominant widow's peak trait, resulting in specific genotype frequencies that reflect the genetic makeup of the population.
\item \textbf{Answer:} Transgressive segregation is a genetic phenomenon that occurs when the offspring of two distinct parental lines exhibit traits that are more extreme than those found in either parent. In the context of the F$_{2}$ generation, this process can lead to a wider range of phenotypic expressions, resulting in offspring that may display traits beyond the parental extremes. However, contrary to the provided key concept that states there is no variability in the F$_{2}$ generation, transgressive segregation actually enhances genetic variability by allowing for the combination of alleles in new ways, thus producing individuals with novel traits. This increased variability has important implications for evolution and adaptation, as it may provide a greater pool of traits for natural selection to act upon. Overall, transgressive segregation underscores the complexity of inheritance and the potential for diversity within populations.
\\ \textit{Note:} correct because transgressive segregation results in the F$_{2}$ generation exhibiting a broader range of phenotypic traits, including extremes that surpass those of the parental lines, thus enhancing genetic variability rather than limiting it.
\end{enumerate}

\vfill
\noindent\textit{End of Answer Key}
\end{document}